% Options for packages loaded elsewhere
\PassOptionsToPackage{unicode}{hyperref}
\PassOptionsToPackage{hyphens}{url}
\documentclass[
  11pt,
  a4paper,
]{article}
\usepackage{xcolor}
\usepackage[margin=18mm]{geometry}
\usepackage{amsmath,amssymb}
\setcounter{secnumdepth}{-\maxdimen} % remove section numbering
\usepackage{iftex}
\ifPDFTeX
  \usepackage[T1]{fontenc}
  \usepackage[utf8]{inputenc}
  \usepackage{textcomp} % provide euro and other symbols
\else % if luatex or xetex
  \usepackage{unicode-math} % this also loads fontspec
  \defaultfontfeatures{Scale=MatchLowercase}
  \defaultfontfeatures[\rmfamily]{Ligatures=TeX,Scale=1}
\fi
\ifPDFTeX\else
  % xetex/luatex font selection
    \setmainfont[]{Noto Serif}
    \setsansfont[]{Noto Sans}
    \setmonofont[]{Noto Sans Mono}
  \setmathfont[]{Noto Sans Math}
\fi
% Use upquote if available, for straight quotes in verbatim environments
\IfFileExists{upquote.sty}{\usepackage{upquote}}{}
\IfFileExists{microtype.sty}{% use microtype if available
  \usepackage[]{microtype}
  \UseMicrotypeSet[protrusion]{basicmath} % disable protrusion for tt fonts
}{}
\makeatletter
\@ifundefined{KOMAClassName}{% if non-KOMA class
  \IfFileExists{parskip.sty}{%
    \usepackage{parskip}
  }{% else
    \setlength{\parindent}{0pt}
    \setlength{\parskip}{6pt plus 2pt minus 1pt}}
}{% if KOMA class
  \KOMAoptions{parskip=half}}
\makeatother
\ifLuaTeX
\usepackage[bidi=basic]{babel}
\else
\usepackage[bidi=default]{babel}
\fi
\babelprovide[main,import]{russian}
\ifPDFTeX
\else
\babelfont{rm}[]{Noto Serif}
\fi
% get rid of language-specific shorthands (see #6817):
\let\LanguageShortHands\languageshorthands
\def\languageshorthands#1{}
\setlength{\emergencystretch}{3em} % prevent overfull lines
\providecommand{\tightlist}{%
  \setlength{\itemsep}{0pt}\setlength{\parskip}{0pt}}
\usepackage{bookmark}
\IfFileExists{xurl.sty}{\usepackage{xurl}}{} % add URL line breaks if available
\urlstyle{same}
\hypersetup{
  pdftitle={Билеты по дискретным случайным процессам},
  pdflang={ru-RU},
  hidelinks,
  pdfcreator={LaTeX via pandoc}}

\title{Билеты по дискретным случайным процессам}
\author{}
\date{}

\begin{document}
\maketitle

\section{Билет 23. Марковские случайные цепи и их
характеристики}\label{ux431ux438ux43bux435ux442-23.-ux43cux430ux440ux43aux43eux432ux441ux43aux438ux435-ux441ux43bux443ux447ux430ux439ux43dux44bux435-ux446ux435ux43fux438-ux438-ux438ux445-ux445ux430ux440ux430ux43aux442ux435ux440ux438ux441ux442ux438ux43aux438}

Источники: Ширяев 2, гл. VIII, §1, §4-7; Сердобольская, лекция 5.

\subsection{1. Короткий
ответ}\label{ux43aux43eux440ux43eux442ux43aux438ux439-ux43eux442ux432ux435ux442}

Марковская цепь - это случайная последовательность

\[
X_0,X_1,X_2,\ldots
\]

со значениями в конечном или счетном пространстве состояний \(E\), у
которой условное распределение будущего при известном настоящем не
зависит от прошлого:

\[
P\{X_{n+1}=j\mid X_0=i_0,\ldots,X_n=i\}
=
P\{X_{n+1}=j\mid X_n=i\}.
\]

Если цепь однородна по времени, то

\[
P\{X_{n+1}=j\mid X_n=i\}=p_{ij},
\]

где

\[
P=(p_{ij})
\]

\begin{itemize}
\tightlist
\item
  стохастическая матрица переходов:
\end{itemize}

\[
p_{ij}\ge0,
\qquad
\sum_jp_{ij}=1.
\]

Закон цепи задается начальным распределением \(\pi_0\) и матрицей \(P\).
Для траектории:

\[
P\{X_0=i_0,\ldots,X_n=i_n\}
=
\pi_0(i_0)p_{i_0i_1}\cdots p_{i_{n-1}i_n}.
\]

Распределения развиваются по формуле

\[
\pi_{n+1}=\pi_nP,
\qquad
\pi_n=\pi_0P^n.
\]

Главные характеристики состояний: достижимость, сообщаемость, классы
сообщающихся состояний, неразложимость, период,
возвратность/невозвратность, положительная возвратность. Стационарное
распределение \(\pi\) удовлетворяет

\[
\pi=\pi P.
\]

Для конечной неразложимой цепи существует единственное стационарное
распределение \(\pi\). Если цепь дополнительно апериодична, то

\[
p_{ij}^{(n)}\to\pi_j.
\]

\subsection{2. Полный
ответ}\label{ux43fux43eux43bux43dux44bux439-ux43eux442ux432ux435ux442}

\subsubsection{2.1. Понятие марковской цепи и марковское
свойство}\label{ux43fux43eux43dux44fux442ux438ux435-ux43cux430ux440ux43aux43eux432ux441ux43aux43eux439-ux446ux435ux43fux438-ux438-ux43cux430ux440ux43aux43eux432ux441ux43aux43eux435-ux441ux432ux43eux439ux441ux442ux432ux43e}

Марковские цепи - основной класс случайных последовательностей с
зависимостью. В iid-последовательности все наблюдения независимы. В
марковской цепи зависимость есть, но она локальная: для прогноза
следующего состояния достаточно знать текущее состояние.

Пусть пространство состояний

\[
E=\{1,2,\ldots\}
\]

конечно или счетно. Последовательность случайных величин

\[
X=(X_n)_{n\ge0}
\]

со значениями в \(E\) называется марковской цепью, если для любых
допустимых состояний

\[
i_0,\ldots,i_n,j
\]

при условии, что условная вероятность определена, выполняется марковское
свойство:

\[
P\{X_{n+1}=j\mid X_0=i_0,\ldots,X_{n-1}=i_{n-1},X_n=i_n\}
=
P\{X_{n+1}=j\mid X_n=i_n\}.
\]

Интуиция: прошлое влияет на будущее только через настоящее состояние.
Если известно \(X_n\), то дополнительная информация о
\(X_0,\ldots,X_{n-1}\) не меняет условный закон \(X_{n+1}\).

\subsubsection{2.2. Однородность и матрица переходных
вероятностей}\label{ux43eux434ux43dux43eux440ux43eux434ux43dux43eux441ux442ux44c-ux438-ux43cux430ux442ux440ux438ux446ux430-ux43fux435ux440ux435ux445ux43eux434ux43dux44bux445-ux432ux435ux440ux43eux44fux442ux43dux43eux441ux442ux435ux439}

Цепь называется однородной по времени, если эти условные вероятности не
зависят от \(n\):

\[
P\{X_{n+1}=j\mid X_n=i\}=p_{ij}.
\]

Матрица

\[
P=(p_{ij})_{i,j\in E}
\]

называется матрицей переходных вероятностей. Она уже подробно
обсуждалась в Билете 22. Здесь важно помнить:

\[
p_{ij}\ge0,
\qquad
\sum_jp_{ij}=1.
\]

В строковом соглашении распределение в момент \(n\) записывается строкой

\[
\pi_n=(\pi_n(i))_{i\in E},
\qquad
\pi_n(i)=P\{X_n=i\}.
\]

Тогда

\[
\pi_{n+1}=\pi_nP,
\qquad
\pi_n=\pi_0P^n.
\]

Покомпонентно:

\[
\pi_{n+1}(j)=\sum_i\pi_n(i)p_{ij}.
\]

\subsubsection{2.3. Описание цепи и формула распределения
траекторий}\label{ux43eux43fux438ux441ux430ux43dux438ux435-ux446ux435ux43fux438-ux438-ux444ux43eux440ux43cux443ux43bux430-ux440ux430ux441ux43fux440ux435ux434ux435ux43bux435ux43dux438ux44f-ux442ux440ux430ux435ux43aux442ux43eux440ux438ux439}

Марковская цепь полностью задается двумя объектами:

\begin{enumerate}
\def\labelenumi{\arabic{enumi}.}
\tightlist
\item
  начальным распределением
\end{enumerate}

\[
\pi_0(i)=P\{X_0=i\};
\]

\begin{enumerate}
\def\labelenumi{\arabic{enumi}.}
\setcounter{enumi}{1}
\tightlist
\item
  матрицей переходов \(P=(p_{ij})\).
\end{enumerate}

Действительно, для любой конечной траектории

\[
i_0,i_1,\ldots,i_n
\]

имеет место формула:

\[
P\{X_0=i_0,\ldots,X_n=i_n\}
=
\pi_0(i_0)p_{i_0i_1}p_{i_1i_2}\cdots p_{i_{n-1}i_n}.
\]

Это ключевая формула билета. Она получается последовательным применением
правила умножения вероятностей и марковского свойства:

\[
P\{X_k=i_k\mid X_0=i_0,\ldots,X_{k-1}=i_{k-1}\}
=
p_{i_{k-1}i_k}.
\]

Важно не путать марковость с независимостью. В марковской цепи состояний
обычно зависимы. Например, если \(p_{ii}\) велико, то знание \(X_n=i\)
сильно влияет на распределение \(X_{n+1}\). Марковость говорит не об
отсутствии зависимости, а о специальной форме зависимости: будущее
условно независимо от прошлого при фиксированном настоящем.

\subsubsection{2.4. Переходные вероятности за n шагов и уравнение
Чепмена-Колмогорова}\label{ux43fux435ux440ux435ux445ux43eux434ux43dux44bux435-ux432ux435ux440ux43eux44fux442ux43dux43eux441ux442ux438-ux437ux430-n-ux448ux430ux433ux43eux432-ux438-ux443ux440ux430ux432ux43dux435ux43dux438ux435-ux447ux435ux43fux43cux435ux43dux430-ux43aux43eux43bux43cux43eux433ux43eux440ux43eux432ux430}

Для анализа цепи изучают не только одношаговые переходы \(p_{ij}\), но и
\(n\)-шаговые переходы:

\[
p_{ij}^{(n)}=P\{X_n=j\mid X_0=i\}.
\]

Для однородной цепи

\[
P^{(n)}=P^n,
\]

и выполнено уравнение Чепмена-Колмогорова:

\[
p_{ij}^{(m+n)}
=
\sum_kp_{ik}^{(n)}p_{kj}^{(m)}.
\]

\subsubsection{2.5. Классификация состояний: сообщаемость и
неразложимость}\label{ux43aux43bux430ux441ux441ux438ux444ux438ux43aux430ux446ux438ux44f-ux441ux43eux441ux442ux43eux44fux43dux438ux439-ux441ux43eux43eux431ux449ux430ux435ux43cux43eux441ux442ux44c-ux438-ux43dux435ux440ux430ux437ux43bux43eux436ux438ux43cux43eux441ux442ux44c}

Теперь перейдем к характеристикам состояний.

Состояние \(j\) называется достижимым из состояния \(i\), если
существует \(n\ge0\) такое, что

\[
p_{ij}^{(n)}>0.
\]

Пишут

\[
i\to j.
\]

Состояния \(i\) и \(j\) называются сообщающимися, если

\[
i\to j
\quad\text{и}\quad
j\to i.
\]

Пишут

\[
i\leftrightarrow j.
\]

Сообщаемость является отношением эквивалентности: она рефлексивна,
симметрична и транзитивна. Поэтому пространство состояний разбивается на
классы сообщающихся состояний.

Цепь называется неразложимой, если любые два состояния сообщаются:

\[
\forall i,j\in E
\quad
i\leftrightarrow j.
\]

Иными словами, из любого состояния можно попасть в любое другое с
положительной вероятностью за некоторое число шагов.

\subsubsection{2.6. Замкнутые классы и периодичность
состояний}\label{ux437ux430ux43cux43aux43dux443ux442ux44bux435-ux43aux43bux430ux441ux441ux44b-ux438-ux43fux435ux440ux438ux43eux434ux438ux447ux43dux43eux441ux442ux44c-ux441ux43eux441ux442ux43eux44fux43dux438ux439}

Класс состояний называется замкнутым, если из него нельзя выйти:

\[
i\in C,\ i\to j
\quad\Rightarrow\quad
j\in C.
\]

Если цепь попала в замкнутый класс, она остается в нем навсегда.
Одноточечный замкнутый класс называется поглощающим состоянием:

\[
p_{ii}=1.
\]

Еще одна важная характеристика - период. Период состояния \(i\):

\[
d(i)=\gcd\{n\ge1:p_{ii}^{(n)}>0\}.
\]

Если

\[
d(i)=1,
\]

состояние называется апериодическим. Если \(d(i)>1\), возвращения в
\(i\) возможны только в моменты времени, кратные некоторому числу или
лежащие в одном циклическом классе. В одном классе сообщающихся
состояний все состояния имеют один и тот же период.

\subsubsection{2.7. Возвратность и невозвратность: определения и
критерии}\label{ux432ux43eux437ux432ux440ux430ux442ux43dux43eux441ux442ux44c-ux438-ux43dux435ux432ux43eux437ux432ux440ux430ux442ux43dux43eux441ux442ux44c-ux43eux43fux440ux435ux434ux435ux43bux435ux43dux438ux44f-ux438-ux43aux440ux438ux442ux435ux440ux438ux438}

Теперь возвратность. Пусть

\[
\tau_i=\inf\{n\ge1:X_n=i\}
\]

\begin{itemize}
\tightlist
\item
  время первого возвращения в \(i\) после выхода из \(i\); если
  возвращения нет, считают \(\tau_i=\infty\). Если цепь стартует из
  \(i\), то вероятность когда-нибудь вернуться в \(i\) равна
\end{itemize}

\[
f_{ii}=P_i\{\tau_i<\infty\}.
\]

Состояние \(i\) называется возвратным, если

\[
f_{ii}=1.
\]

Состояние \(i\) называется невозвратным, или транзиентным, если

\[
f_{ii}<1.
\]

Интуитивно: возвратное состояние почти наверное посещается снова, если
из него стартовать. Невозвратное состояние с положительной вероятностью
покидается навсегда.

Критерий через \(n\)-шаговые вероятности:

\[
i\ \text{возвратно}
\quad\Longleftrightarrow\quad
\sum_{n=0}^{\infty}p_{ii}^{(n)}=\infty.
\]

И

\[
i\ \text{невозвратно}
\quad\Longleftrightarrow\quad
\sum_{n=0}^{\infty}p_{ii}^{(n)}<\infty.
\]

\subsubsection{2.8. Среднее время возвращения и типы
возвратности}\label{ux441ux440ux435ux434ux43dux435ux435-ux432ux440ux435ux43cux44f-ux432ux43eux437ux432ux440ux430ux449ux435ux43dux438ux44f-ux438-ux442ux438ux43fux44b-ux432ux43eux437ux432ux440ux430ux442ux43dux43eux441ux442ux438}

Для возвратного состояния вводят среднее время возвращения:

\[
\mu_i=E_i\tau_i.
\]

Если

\[
\mu_i<\infty,
\]

то состояние называется положительно возвратным. Если

\[
\mu_i=\infty,
\]

то оно называется нулевым возвратным.

Для конечных неразложимых цепей все состояния положительно возвратны.
Для счетных цепей это уже не всегда верно: например, симметричное
случайное блуждание на \(\mathbb Z\) возвратно, но не положительно
возвратно.

\subsubsection{2.9. Стационарное распределение и равновесный
режим}\label{ux441ux442ux430ux446ux438ux43eux43dux430ux440ux43dux43eux435-ux440ux430ux441ux43fux440ux435ux434ux435ux43bux435ux43dux438ux435-ux438-ux440ux430ux432ux43dux43eux432ux435ux441ux43dux44bux439-ux440ux435ux436ux438ux43c}

Стационарное распределение - это вероятностный вектор

\[
\pi=(\pi_i)_{i\in E}
\]

такой, что

\[
\pi_i\ge0,
\qquad
\sum_i\pi_i=1,
\]

и

\[
\pi=\pi P.
\]

Покомпонентно:

\[
\pi_j=\sum_i\pi_i p_{ij}.
\]

Если \(X_0\) имеет распределение \(\pi\), то

\[
\pi_1=\pi P=\pi.
\]

Значит все \(X_n\) имеют одно и то же распределение \(\pi\). Поэтому
стационарное распределение описывает равновесный режим цепи. Важно:
стационарное распределение цепи - это не то же самое, что стационарность
любой конкретной реализации. Это распределение, которое сохраняется
переходом.

\subsubsection{2.10. Предельное поведение и условия
эргодичности}\label{ux43fux440ux435ux434ux435ux43bux44cux43dux43eux435-ux43fux43eux432ux435ux434ux435ux43dux438ux435-ux438-ux443ux441ux43bux43eux432ux438ux44f-ux44dux440ux433ux43eux434ux438ux447ux43dux43eux441ux442ux438}

Основная предельная картина такая:

\begin{itemize}
\tightlist
\item
  если цепь конечна, неразложима и апериодична, то существует
  единственное стационарное распределение \(\pi\);
\item
  для любых начальных состояний \(i\) и любых состояний \(j\)
\end{itemize}

\[
p_{ij}^{(n)}\to\pi_j;
\]

\begin{itemize}
\tightlist
\item
  для любого начального распределения \(\pi_0\)
\end{itemize}

\[
\pi_0P^n\to\pi.
\]

В счетном случае нужны более тонкие условия. По Ширяеву: существование
единственного стационарного распределения связано с наличием
единственного неразложимого положительно возвратного класса, а
существование эргодического предельного распределения требует
неразложимости, положительной возвратности и апериодичности.

\subsection{3. Основные
определения}\label{ux43eux441ux43dux43eux432ux43dux44bux435-ux43eux43fux440ux435ux434ux435ux43bux435ux43dux438ux44f}

\subsubsection{Марковская
цепь}\label{ux43cux430ux440ux43aux43eux432ux441ux43aux430ux44f-ux446ux435ux43fux44c}

Последовательность \((X_n)_{n\ge0}\) со значениями в \(E\) называется
марковской цепью, если

\[
P\{X_{n+1}=j\mid X_0=i_0,\ldots,X_n=i_n\}
=
P\{X_{n+1}=j\mid X_n=i_n\}.
\]

\subsubsection{Однородная марковская
цепь}\label{ux43eux434ux43dux43eux440ux43eux434ux43dux430ux44f-ux43cux430ux440ux43aux43eux432ux441ux43aux430ux44f-ux446ux435ux43fux44c}

Цепь называется однородной, если

\[
P\{X_{n+1}=j\mid X_n=i\}=p_{ij}
\]

не зависит от \(n\).

\subsubsection{Матрица
переходов}\label{ux43cux430ux442ux440ux438ux446ux430-ux43fux435ux440ux435ux445ux43eux434ux43eux432}

Матрица

\[
P=(p_{ij}),
\qquad
p_{ij}=P\{X_{n+1}=j\mid X_n=i\}.
\]

Она стохастическая по строкам:

\[
p_{ij}\ge0,
\qquad
\sum_jp_{ij}=1.
\]

\subsubsection{Начальное
распределение}\label{ux43dux430ux447ux430ux43bux44cux43dux43eux435-ux440ux430ux441ux43fux440ux435ux434ux435ux43bux435ux43dux438ux435}

Вероятностный вектор

\[
\pi_0(i)=P\{X_0=i\}.
\]

\subsubsection{\texorpdfstring{Распределение в момент
\(n\)}{Распределение в момент n}}\label{ux440ux430ux441ux43fux440ux435ux434ux435ux43bux435ux43dux438ux435-ux432-ux43cux43eux43cux435ux43dux442-n}

\[
\pi_n(i)=P\{X_n=i\}.
\]

Для строкового соглашения:

\[
\pi_n=\pi_0P^n.
\]

\subsubsection{\texorpdfstring{\(n\)-шаговая переходная
вероятность}{n-шаговая переходная вероятность}}\label{n-ux448ux430ux433ux43eux432ux430ux44f-ux43fux435ux440ux435ux445ux43eux434ux43dux430ux44f-ux432ux435ux440ux43eux44fux442ux43dux43eux441ux442ux44c}

\[
p_{ij}^{(n)}=P\{X_n=j\mid X_0=i\}.
\]

\subsubsection{Достижимость}\label{ux434ux43eux441ux442ux438ux436ux438ux43cux43eux441ux442ux44c}

Состояние \(j\) достижимо из \(i\), если

\[
\exists n\ge0:\ p_{ij}^{(n)}>0.
\]

Обозначение:

\[
i\to j.
\]

\subsubsection{Сообщаемость}\label{ux441ux43eux43eux431ux449ux430ux435ux43cux43eux441ux442ux44c}

Состояния \(i\) и \(j\) сообщаются, если

\[
i\to j
\quad\text{и}\quad
j\to i.
\]

Обозначение:

\[
i\leftrightarrow j.
\]

\subsubsection{Класс сообщающихся
состояний}\label{ux43aux43bux430ux441ux441-ux441ux43eux43eux431ux449ux430ux44eux449ux438ux445ux441ux44f-ux441ux43eux441ux442ux43eux44fux43dux438ux439}

Класс эквивалентности по отношению сообщаемости.

\subsubsection{Неразложимая
цепь}\label{ux43dux435ux440ux430ux437ux43bux43eux436ux438ux43cux430ux44f-ux446ux435ux43fux44c}

Цепь называется неразложимой, если все состояния принадлежат одному
классу сообщаемости:

\[
\forall i,j\in E:\ i\leftrightarrow j.
\]

\subsubsection{Замкнутый
класс}\label{ux437ux430ux43cux43aux43dux443ux442ux44bux439-ux43aux43bux430ux441ux441}

Класс \(C\) замкнут, если из него нельзя выйти:

\[
i\in C,\ p_{ij}^{(n)}>0
\quad\Rightarrow\quad
j\in C.
\]

Для одношаговых переходов достаточно проверять:

\[
i\in C,\ p_{ij}>0
\quad\Rightarrow\quad
j\in C.
\]

\subsubsection{Поглощающее
состояние}\label{ux43fux43eux433ux43bux43eux449ux430ux44eux449ux435ux435-ux441ux43eux441ux442ux43eux44fux43dux438ux435}

Состояние \(i\) поглощающее, если

\[
p_{ii}=1.
\]

\subsubsection{Период
состояния}\label{ux43fux435ux440ux438ux43eux434-ux441ux43eux441ux442ux43eux44fux43dux438ux44f}

\[
d(i)=\gcd\{n\ge1:p_{ii}^{(n)}>0\}.
\]

Если \(d(i)=1\), состояние апериодично. Если \(d(i)>1\), состояние
периодично.

\subsubsection{Время первого
возвращения}\label{ux432ux440ux435ux43cux44f-ux43fux435ux440ux432ux43eux433ux43e-ux432ux43eux437ux432ux440ux430ux449ux435ux43dux438ux44f}

При старте из \(i\):

\[
\tau_i=\inf\{n\ge1:X_n=i\}.
\]

\subsubsection{Возвратное
состояние}\label{ux432ux43eux437ux432ux440ux430ux442ux43dux43eux435-ux441ux43eux441ux442ux43eux44fux43dux438ux435}

Состояние \(i\) возвратно, если

\[
P_i\{\tau_i<\infty\}=1.
\]

\subsubsection{Невозвратное
состояние}\label{ux43dux435ux432ux43eux437ux432ux440ux430ux442ux43dux43eux435-ux441ux43eux441ux442ux43eux44fux43dux438ux435}

Состояние \(i\) невозвратно, если

\[
P_i\{\tau_i<\infty\}<1.
\]

\subsubsection{Положительно возвратное
состояние}\label{ux43fux43eux43bux43eux436ux438ux442ux435ux43bux44cux43dux43e-ux432ux43eux437ux432ux440ux430ux442ux43dux43eux435-ux441ux43eux441ux442ux43eux44fux43dux438ux435}

Возвратное состояние \(i\) положительно возвратно, если

\[
E_i\tau_i<\infty.
\]

\subsubsection{Нулевое возвратное
состояние}\label{ux43dux443ux43bux435ux432ux43eux435-ux432ux43eux437ux432ux440ux430ux442ux43dux43eux435-ux441ux43eux441ux442ux43eux44fux43dux438ux435}

Возвратное состояние \(i\) нулевое возвратное, если

\[
E_i\tau_i=\infty.
\]

\subsubsection{Стационарное
распределение}\label{ux441ux442ux430ux446ux438ux43eux43dux430ux440ux43dux43eux435-ux440ux430ux441ux43fux440ux435ux434ux435ux43bux435ux43dux438ux435}

Вероятностный вектор \(\pi\) такой, что

\[
\pi=\pi P.
\]

Покомпонентно:

\[
\pi_j=\sum_i\pi_i p_{ij}.
\]

\subsubsection{Эргодическая цепь в конечном
случае}\label{ux44dux440ux433ux43eux434ux438ux447ux435ux441ux43aux430ux44f-ux446ux435ux43fux44c-ux432-ux43aux43eux43dux435ux447ux43dux43eux43c-ux441ux43bux443ux447ux430ux435}

Часто конечную цепь называют эргодической, если она неразложима и
апериодична. Тогда существует единственное стационарное распределение и

\[
p_{ij}^{(n)}\to\pi_j.
\]

\subsection{4. Основные теоремы и
свойства}\label{ux43eux441ux43dux43eux432ux43dux44bux435-ux442ux435ux43eux440ux435ux43cux44b-ux438-ux441ux432ux43eux439ux441ux442ux432ux430}

\subsubsection{1. Закон
траектории}\label{ux437ux430ux43aux43eux43d-ux442ux440ux430ux435ux43aux442ux43eux440ux438ux438}

Если цепь однородна, имеет начальное распределение \(\pi_0\) и матрицу
переходов \(P\), то

\[
P\{X_0=i_0,\ldots,X_n=i_n\}
=
\pi_0(i_0)\prod_{k=1}^n p_{i_{k-1}i_k}.
\]

\subsubsection{2. Эволюция
распределений}\label{ux44dux432ux43eux43bux44eux446ux438ux44f-ux440ux430ux441ux43fux440ux435ux434ux435ux43bux435ux43dux438ux439}

Для строковых распределений:

\[
\pi_{n+1}=\pi_nP,
\qquad
\pi_n=\pi_0P^n.
\]

\subsubsection{3.
Чепмен-Колмогоров}\label{ux447ux435ux43fux43cux435ux43d-ux43aux43eux43bux43cux43eux433ux43eux440ux43eux432}

Для \(n\)-шаговых переходов:

\[
p_{ij}^{(m+n)}
=
\sum_kp_{ik}^{(n)}p_{kj}^{(m)}.
\]

Эта формула доказывает, что

\[
P^{(n)}=P^n.
\]

\subsubsection{4. Сообщаемость - отношение
эквивалентности}\label{ux441ux43eux43eux431ux449ux430ux435ux43cux43eux441ux442ux44c---ux43eux442ux43dux43eux448ux435ux43dux438ux435-ux44dux43aux432ux438ux432ux430ux43bux435ux43dux442ux43dux43eux441ux442ux438}

Свойство

\[
i\leftrightarrow j
\]

рефлексивно, симметрично и транзитивно. Транзитивность следует из
Чепмена-Колмогорова: если \(p_{ij}^{(n)}>0\) и \(p_{jk}^{(m)}>0\), то

\[
p_{ik}^{(n+m)}
\ge
p_{ij}^{(n)}p_{jk}^{(m)}
>0.
\]

\subsubsection{5. В одном классе сообщаемости свойства
совпадают}\label{ux432-ux43eux434ux43dux43eux43c-ux43aux43bux430ux441ux441ux435-ux441ux43eux43eux431ux449ux430ux435ux43cux43eux441ux442ux438-ux441ux432ux43eux439ux441ux442ux432ux430-ux441ux43eux432ux43fux430ux434ux430ux44eux442}

Если \(i\leftrightarrow j\), то состояния \(i\) и \(j\) имеют один и тот
же период. Также возвратность и невозвратность являются свойствами
класса: в одном классе либо все состояния возвратны, либо все
невозвратны.

\subsubsection{6. Критерий
возвратности}\label{ux43aux440ux438ux442ux435ux440ux438ux439-ux432ux43eux437ux432ux440ux430ux442ux43dux43eux441ux442ux438}

Состояние \(i\) возвратно тогда и только тогда, когда

\[
\sum_{n=0}^{\infty}p_{ii}^{(n)}=\infty.
\]

Состояние \(i\) невозвратно тогда и только тогда, когда

\[
\sum_{n=0}^{\infty}p_{ii}^{(n)}<\infty.
\]

\subsubsection{7. Стационарное распределение сохраняется
переходом}\label{ux441ux442ux430ux446ux438ux43eux43dux430ux440ux43dux43eux435-ux440ux430ux441ux43fux440ux435ux434ux435ux43bux435ux43dux438ux435-ux441ux43eux445ux440ux430ux43dux44fux435ux442ux441ux44f-ux43fux435ux440ux435ux445ux43eux434ux43eux43c}

Если

\[
\pi=\pi P,
\]

и

\[
\pi_0=\pi,
\]

то

\[
\pi_n=\pi
\]

для всех \(n\).

\subsubsection{8. Конечная неразложимая
цепь}\label{ux43aux43eux43dux435ux447ux43dux430ux44f-ux43dux435ux440ux430ux437ux43bux43eux436ux438ux43cux430ux44f-ux446ux435ux43fux44c}

Если \(E\) конечно и цепь неразложима, то существует единственное
стационарное распределение \(\pi\), причем

\[
\pi_i>0
\]

для всех \(i\). Все состояния положительно возвратны.

\subsubsection{9. Конечная неразложимая апериодическая
цепь}\label{ux43aux43eux43dux435ux447ux43dux430ux44f-ux43dux435ux440ux430ux437ux43bux43eux436ux438ux43cux430ux44f-ux430ux43fux435ux440ux438ux43eux434ux438ux447ux435ux441ux43aux430ux44f-ux446ux435ux43fux44c}

Если \(E\) конечно, цепь неразложима и апериодична, то

\[
p_{ij}^{(n)}\to\pi_j
\]

для всех \(i,j\). То есть предельное распределение не зависит от
начального состояния.

\subsubsection{10. Счетный
случай}\label{ux441ux447ux435ux442ux43dux44bux439-ux441ux43bux443ux447ux430ux439}

Для счетной цепи существование единственного стационарного распределения
требует положительной возвратности соответствующего замкнутого
неразложимого класса. Для эргодического предела дополнительно нужна
апериодичность.

\subsection{5. Примеры}\label{ux43fux440ux438ux43cux435ux440ux44b}

\subsubsection{Пример 1. Двухсостоянийная
цепь}\label{ux43fux440ux438ux43cux435ux440-1.-ux434ux432ux443ux445ux441ux43eux441ux442ux43eux44fux43dux438ux439ux43dux430ux44f-ux446ux435ux43fux44c}

Пусть

\[
E=\{0,1\},
\]

и

\[
P=
\begin{pmatrix}
1-a & a\\
b & 1-b
\end{pmatrix},
\qquad
0<a,b<1.
\]

Из любого состояния можно попасть в любое другое:

\[
0\leftrightarrow1.
\]

Цепь неразложима. Так как

\[
p_{00}=1-a>0,
\qquad
p_{11}=1-b>0,
\]

то возврат в каждое состояние возможен за один шаг, и период равен
\(1\). Цепь апериодична.

Стационарное распределение \(\pi=(\pi_0,\pi_1)\) находится из

\[
\pi=\pi P,
\qquad
\pi_0+\pi_1=1.
\]

Получаем

\[
\pi_0a=\pi_1b,
\]

поэтому

\[
\pi_0=\frac{b}{a+b},
\qquad
\pi_1=\frac{a}{a+b}.
\]

\subsubsection{Пример 2. Поглощающее
состояние}\label{ux43fux440ux438ux43cux435ux440-2.-ux43fux43eux433ux43bux43eux449ux430ux44eux449ux435ux435-ux441ux43eux441ux442ux43eux44fux43dux438ux435}

Пусть

\[
P=
\begin{pmatrix}
1 & 0\\
q & 1-q
\end{pmatrix},
\qquad
0<q\le1.
\]

Состояние \(0\) поглощающее:

\[
p_{00}=1.
\]

Из состояния \(1\) можно попасть в \(0\), но из \(0\) нельзя вернуться в
\(1\). Поэтому цепь разложима. Стационарное распределение:

\[
\pi=(1,0).
\]

Состояние \(1\) невозвратно: если цепь ушла из \(1\) в \(0\), она уже не
вернется.

\subsubsection{\texorpdfstring{Пример 3. Простое случайное блуждание на
\(\mathbb Z\)}{Пример 3. Простое случайное блуждание на \textbackslash mathbb Z}}\label{ux43fux440ux438ux43cux435ux440-3.-ux43fux440ux43eux441ux442ux43eux435-ux441ux43bux443ux447ux430ux439ux43dux43eux435-ux431ux43bux443ux436ux434ux430ux43dux438ux435-ux43dux430-mathbb-z}

Пусть

\[
p_{i,i+1}=p,
\qquad
p_{i,i-1}=q=1-p.
\]

Если \(0<p<1\), все состояния сообщаются, цепь неразложима. Возвраты в
исходное состояние возможны только за четное число шагов, поэтому период
равен

\[
d=2.
\]

При \(p=q=1/2\) блуждание на \(\mathbb Z\) возвратно, но среднее время
возвращения бесконечно, поэтому оно нулевое возвратное и не имеет
стационарного вероятностного распределения.

При \(p\ne q\) блуждание на \(\mathbb Z\) невозвратно.

\subsubsection{Пример 4. Случайное блуждание на конечном
графе}\label{ux43fux440ux438ux43cux435ux440-4.-ux441ux43bux443ux447ux430ux439ux43dux43eux435-ux431ux43bux443ux436ux434ux430ux43dux438ux435-ux43dux430-ux43aux43eux43dux435ux447ux43dux43eux43c-ux433ux440ux430ux444ux435}

Пусть есть конечный связный неориентированный граф. Из вершины \(i\)
цепь переходит равновероятно в одну из соседних вершин:

\[
p_{ij}=
\begin{cases}
\frac{1}{\deg(i)}, & i\sim j,\\
0, & \text{иначе}.
\end{cases}
\]

Если граф связный, цепь неразложима. Стационарное распределение:

\[
\pi_i=\frac{\deg(i)}{\sum_k\deg(k)}.
\]

Если граф двудольный, цепь имеет период \(2\). Если в графе есть петля
или нарушена двудольность, часто получается апериодичность.

\subsubsection{Пример 5. Цепь
рождения-гибели}\label{ux43fux440ux438ux43cux435ux440-5.-ux446ux435ux43fux44c-ux440ux43eux436ux434ux435ux43dux438ux44f-ux433ux438ux431ux435ux43bux438}

На состояниях

\[
E=\{0,1,2,\ldots\}
\]

переходы возможны только в соседние состояния:

\[
p_{i,i+1}=p_i,
\qquad
p_{i,i-1}=q_i,
\qquad
p_{ii}=r_i,
\]

где

\[
p_i+q_i+r_i=1.
\]

Такие цепи моделируют численность очереди, популяции или запаса. Их
характеристики зависят от соотношения вероятностей рождения и гибели:
цепь может иметь стационарное распределение, а может уходить на
бесконечность.

\subsection{6. Схемы доказательств /
обоснований}\label{ux441ux445ux435ux43cux44b-ux434ux43eux43aux430ux437ux430ux442ux435ux43bux44cux441ux442ux432-ux43eux431ux43eux441ux43dux43eux432ux430ux43dux438ux439}

\textbf{Доказательство формулы траектории.}

По правилу умножения вероятностей:

\[
P\{X_0=i_0,\ldots,X_n=i_n\}
=
P\{X_0=i_0\}
\prod_{k=1}^n
P\{X_k=i_k\mid X_0=i_0,\ldots,X_{k-1}=i_{k-1}\}.
\]

По марковскому свойству:

\[
P\{X_k=i_k\mid X_0=i_0,\ldots,X_{k-1}=i_{k-1}\}
=
P\{X_k=i_k\mid X_{k-1}=i_{k-1}\}
=
p_{i_{k-1}i_k}.
\]

Значит

\[
P\{X_0=i_0,\ldots,X_n=i_n\}
=
\pi_0(i_0)\prod_{k=1}^n p_{i_{k-1}i_k}.
\]

\textbf{Доказательство \(\pi_{n+1}=\pi_nP\).}

Для каждого \(j\):

\[
\pi_{n+1}(j)
=
P\{X_{n+1}=j\}
=
\sum_iP\{X_n=i,\ X_{n+1}=j\}
\]

\[
=
\sum_iP\{X_n=i\}P\{X_{n+1}=j\mid X_n=i\}
=
\sum_i\pi_n(i)p_{ij}.
\]

Это ровно \(j\)-я координата строки \(\pi_nP\).

\textbf{Доказательство транзитивности сообщаемости.}

Пусть

\[
i\to j,
\qquad
j\to k.
\]

Тогда существуют \(n,m\) такие, что

\[
p_{ij}^{(n)}>0,
\qquad
p_{jk}^{(m)}>0.
\]

По Чепмену-Колмогорову:

\[
p_{ik}^{(n+m)}
=
\sum_l p_{il}^{(n)}p_{lk}^{(m)}
\ge
p_{ij}^{(n)}p_{jk}^{(m)}
>0.
\]

Значит

\[
i\to k.
\]

Аналогично доказывается обратная достижимость, если есть сообщаемость.

\textbf{Идея критерия возвратности.}

Величина

\[
\sum_{n=0}^{\infty}p_{ii}^{(n)}
\]

равна математическому ожиданию числа посещений состояния \(i\) при
старте из \(i\):

\[
E_i\sum_{n=0}^{\infty}\mathbf 1_{\{X_n=i\}}
=
\sum_{n=0}^{\infty}P_i\{X_n=i\}
=
\sum_{n=0}^{\infty}p_{ii}^{(n)}.
\]

Если состояние невозвратно, число возвращений имеет конечное среднее.
Если состояние возвратно, посещений бесконечно много почти наверное, и
сумма расходится.

\textbf{Почему стационарное распределение сохраняется.}

Если

\[
\pi=\pi P,
\]

и

\[
\pi_0=\pi,
\]

то

\[
\pi_1=\pi_0P=\pi P=\pi.
\]

По индукции:

\[
\pi_n=\pi
\]

для всех \(n\).

\textbf{Идея предельной теоремы для конечной эргодической цепи.}

Неразложимость означает, что цепь не распадается на независимые
замкнутые части: из любого состояния можно добраться до любого другого.
Положительная возвратность в конечном случае гарантирует существование
равновесных долей времени. Апериодичность убирает осцилляции по
четным/нечетным моментам. Поэтому матрица \(P^n\) стабилизируется по
строкам:

\[
(P^n)_{ij}\to\pi_j.
\]

Все строки предельной матрицы одинаковы и равны стационарному
распределению.

\subsection{7. Что написать на
доске}\label{ux447ux442ux43e-ux43dux430ux43fux438ux441ux430ux442ux44c-ux43dux430-ux434ux43eux441ux43aux435}

Марковское свойство:

\[
P\{X_{n+1}=j\mid X_0=i_0,\ldots,X_n=i\}
=
P\{X_{n+1}=j\mid X_n=i\}
=p_{ij}.
\]

Матрица переходов:

\[
p_{ij}\ge0,
\qquad
\sum_jp_{ij}=1.
\]

Закон траектории:

\[
P\{X_0=i_0,\ldots,X_n=i_n\}
=
\pi_0(i_0)p_{i_0i_1}\cdots p_{i_{n-1}i_n}.
\]

Эволюция распределений:

\[
\pi_{n+1}=\pi_nP,
\qquad
\pi_n=\pi_0P^n.
\]

Достижимость и сообщаемость:

\[
i\to j
\Longleftrightarrow
\exists n:\ p_{ij}^{(n)}>0,
\qquad
i\leftrightarrow j
\Longleftrightarrow
i\to j,\ j\to i.
\]

Период:

\[
d(i)=\gcd\{n\ge1:p_{ii}^{(n)}>0\}.
\]

Возвратность:

\[
P_i\{\tau_i<\infty\}=1,
\qquad
\tau_i=\inf\{n\ge1:X_n=i\}.
\]

Стационарное распределение:

\[
\pi=\pi P,
\qquad
\sum_i\pi_i=1.
\]

Предельная формула для конечной эргодической цепи:

\[
p_{ij}^{(n)}\to\pi_j.
\]

\subsection{8. Типичные дополнительные
вопросы}\label{ux442ux438ux43fux438ux447ux43dux44bux435-ux434ux43eux43fux43eux43bux43dux438ux442ux435ux43bux44cux43dux44bux435-ux432ux43eux43fux440ux43eux441ux44b}

\textbf{Марковская цепь - это независимая последовательность?}

Нет. Марковская цепь обычно зависима. Марковское свойство означает
только условную независимость будущего от прошлого при фиксированном
настоящем.

\textbf{Что нужно, чтобы задать закон марковской цепи?}

Нужно начальное распределение \(\pi_0\) и переходная матрица \(P\).
Тогда все конечномерные распределения задаются формулой

\[
P\{X_0=i_0,\ldots,X_n=i_n\}
=
\pi_0(i_0)\prod_{k=1}^n p_{i_{k-1}i_k}.
\]

\textbf{Чем отличается стационарное распределение от стационарной цепи?}

Стационарное распределение \(\pi\) удовлетворяет \(\pi=\pi P\). Если
начальное распределение равно \(\pi\), то процесс становится
стационарным в узком смысле по сдвигам времени. Но сама матрица \(P\)
может иметь стационарное распределение, даже если цепь стартовала не из
него.

\textbf{Что такое неразложимость?}

Это возможность добраться из любого состояния в любое другое:

\[
\forall i,j\quad \exists n:\ p_{ij}^{(n)}>0.
\]

\textbf{Что такое период и зачем он нужен?}

Период состояния \(i\):

\[
d(i)=\gcd\{n\ge1:p_{ii}^{(n)}>0\}.
\]

Если период больше \(1\), вероятности переходов могут не сходиться, а
колебаться. Например, простое блуждание на двудольном графе возвращается
в исходную долю только в четные моменты.

\textbf{Что такое возвратность?}

Состояние возвратно, если при старте из него цепь почти наверное
когда-нибудь вернется:

\[
P_i\{\tau_i<\infty\}=1.
\]

\textbf{Что такое положительная возвратность?}

Это возвратность с конечным средним временем возвращения:

\[
E_i\tau_i<\infty.
\]

Именно положительная возвратность связана с существованием стационарного
распределения.

\textbf{Почему для конечной неразложимой цепи есть стационарное
распределение?}

В конечном случае вероятностный симплекс компактен, а переход
\(\pi\mapsto\pi P\) переводит его в себя. При неразложимости
стационарное распределение единственно и имеет все положительные
координаты.

\textbf{Зачем нужна апериодичность в предельной теореме?}

Без апериодичности цепь может колебаться между циклическими подклассами.
Например, если цепь каждый шаг переходит из \(0\) в \(1\) и из \(1\) в
\(0\), то стационарное распределение есть, но \(P^n\) не сходится.

\textbf{Как связана марковская цепь со случайным блужданием?}

Случайное блуждание

\[
S_{n+1}=S_n+\xi_{n+1}
\]

является марковской цепью, потому что следующий шаг независим от
прошлого, а условное распределение \(S_{n+1}\) зависит только от
\(S_n\).

\subsection{9. Типичные
ошибки}\label{ux442ux438ux43fux438ux447ux43dux44bux435-ux43eux448ux438ux431ux43aux438}

\begin{itemize}
\tightlist
\item
  Считать марковость независимостью состояний.
\item
  Забывать начальное распределение \(\pi_0\).
\item
  Путать матрицу переходов \(P\) и распределение \(\pi_n\).
\item
  Путать стационарность цепи и стационарное распределение.
\item
  Писать \(\pi P=\pi\) без условия \(\sum_i\pi_i=1\) и \(\pi_i\ge0\).
\item
  Называть цепь эргодической только из-за существования стационарного
  распределения. Для сходимости \(p_{ij}^{(n)}\to\pi_j\) нужна еще
  апериодичность и отсутствие разложения.
\item
  Забывать, что в счетном случае не всякая неразложимая цепь имеет
  стационарное вероятностное распределение.
\item
  Путать возвратность и положительную возвратность.
\item
  Считать, что если состояние достижимо из другого, то обратная
  достижимость автоматическая. Это неверно.
\item
  Забывать, что период является свойством класса сообщающихся состояний.
\end{itemize}

\subsection{10. Связи с другими
билетами}\label{ux441ux432ux44fux437ux438-ux441-ux434ux440ux443ux433ux438ux43cux438-ux431ux438ux43bux435ux442ux430ux43cux438}

\begin{itemize}
\tightlist
\item
  Билет 10: марковское свойство формулируется через условные вероятности
  и условные распределения.
\item
  Билет 15: начальное распределение и матрица переходов задают
  согласованные конечномерные распределения.
\item
  Билет 16: марковская цепь - частный случай дискретного случайного
  процесса.
\item
  Билет 17: если цепь стартует из стационарного распределения, она
  стационарна в узком смысле.
\item
  Билет 19: случайное блуждание - важный пример марковской цепи.
\item
  Билет 21: эргодические и предельные свойства цепей являются аналогами
  закона больших чисел для зависимых последовательностей.
\item
  Билет 22: переходные функции и матрицы - техническая основа марковских
  цепей.
\item
  Билет 24: эргодичность стационарных последовательностей и
  корреляционные условия связаны с предельным поведением марковских
  цепей.
\end{itemize}

\subsection{11. Минимум на
удовлетворительно}\label{ux43cux438ux43dux438ux43cux443ux43c-ux43dux430-ux443ux434ux43eux432ux43bux435ux442ux432ux43eux440ux438ux442ux435ux43bux44cux43dux43e}

Нужно сказать:

\begin{enumerate}
\def\labelenumi{\arabic{enumi}.}
\tightlist
\item
  Марковское свойство:
\end{enumerate}

\[
P\{X_{n+1}=j\mid X_0=i_0,\ldots,X_n=i\}
=
P\{X_{n+1}=j\mid X_n=i\}.
\]

\begin{enumerate}
\def\labelenumi{\arabic{enumi}.}
\setcounter{enumi}{1}
\tightlist
\item
  В однородном случае:
\end{enumerate}

\[
P\{X_{n+1}=j\mid X_n=i\}=p_{ij}.
\]

\begin{enumerate}
\def\labelenumi{\arabic{enumi}.}
\setcounter{enumi}{2}
\tightlist
\item
  Закон траектории:
\end{enumerate}

\[
P\{X_0=i_0,\ldots,X_n=i_n\}
=
\pi_0(i_0)p_{i_0i_1}\cdots p_{i_{n-1}i_n}.
\]

\begin{enumerate}
\def\labelenumi{\arabic{enumi}.}
\setcounter{enumi}{3}
\tightlist
\item
  Эволюция распределений:
\end{enumerate}

\[
\pi_{n+1}=\pi_nP.
\]

\begin{enumerate}
\def\labelenumi{\arabic{enumi}.}
\setcounter{enumi}{4}
\tightlist
\item
  Стационарное распределение:
\end{enumerate}

\[
\pi=\pi P.
\]

\subsection{12. Минимум на
хорошо}\label{ux43cux438ux43dux438ux43cux443ux43c-ux43dux430-ux445ux43eux440ux43eux448ux43e}

Помимо минимума, нужно:

\begin{itemize}
\tightlist
\item
  объяснить, почему марковость не равна независимости;
\item
  определить достижимость и сообщаемость;
\item
  сказать, что классы сообщаемости разбивают пространство состояний;
\item
  определить неразложимость и период;
\item
  привести пример поглощающей цепи или случайного блуждания;
\item
  назвать условие сходимости для конечной цепи: неразложимость плюс
  апериодичность.
\end{itemize}

\subsection{13. Что добавить для
отлично}\label{ux447ux442ux43e-ux434ux43eux431ux430ux432ux438ux442ux44c-ux434ux43bux44f-ux43eux442ux43bux438ux447ux43dux43e}

Для сильного ответа стоит добавить:

\begin{itemize}
\tightlist
\item
  возвратность и невозвратность через \(P_i\{\tau_i<\infty\}\);
\item
  критерий возвратности:
\end{itemize}

\[
\sum_{n=0}^{\infty}p_{ii}^{(n)}=\infty;
\]

\begin{itemize}
\tightlist
\item
  положительную и нулевую возвратность;
\item
  точную роль стационарного распределения;
\item
  формулировку предельной теоремы:
\end{itemize}

\[
p_{ij}^{(n)}\to\pi_j;
\]

\begin{itemize}
\tightlist
\item
  предупреждение, что в счетном случае неразложимости и апериодичности
  недостаточно без положительной возвратности.
\end{itemize}

\subsection{14. Устная версия ответа на 5
минут}\label{ux443ux441ux442ux43dux430ux44f-ux432ux435ux440ux441ux438ux44f-ux43eux442ux432ux435ux442ux430-ux43dux430-5-ux43cux438ux43dux443ux442}

Марковская цепь - это случайная последовательность \(X_0,X_1,\ldots\) со
значениями в конечном или счетном множестве состояний \(E\), для которой
будущее при известном настоящем не зависит от прошлого:

\[
P\{X_{n+1}=j\mid X_0=i_0,\ldots,X_n=i\}
=
P\{X_{n+1}=j\mid X_n=i\}.
\]

Если цепь однородна, то правая часть не зависит от \(n\):

\[
P\{X_{n+1}=j\mid X_n=i\}=p_{ij}.
\]

Матрица \(P=(p_{ij})\) стохастическая:

\[
p_{ij}\ge0,
\qquad
\sum_jp_{ij}=1.
\]

Цепь задается начальным распределением \(\pi_0\) и матрицей \(P\). Для
любой траектории:

\[
P\{X_0=i_0,\ldots,X_n=i_n\}
=
\pi_0(i_0)p_{i_0i_1}\cdots p_{i_{n-1}i_n}.
\]

Распределение в момент \(n\) считается по формуле

\[
\pi_n=\pi_0P^n.
\]

Дальше основные характеристики. Состояние \(j\) достижимо из \(i\), если
\(p_{ij}^{(n)}>0\) для некоторого \(n\). Если \(i\) достижимо из \(j\) и
\(j\) достижимо из \(i\), состояния сообщаются. Сообщаемость разбивает
пространство состояний на классы. Если все состояния сообщаются, цепь
неразложима.

Период состояния:

\[
d(i)=\gcd\{n\ge1:p_{ii}^{(n)}>0\}.
\]

Если \(d(i)=1\), состояние апериодично. Период важен для сходимости: при
периоде \(2\) цепь может бесконечно колебаться между двумя частями
пространства.

Возвратность: при старте из \(i\) вводим

\[
\tau_i=\inf\{n\ge1:X_n=i\}.
\]

Состояние возвратно, если

\[
P_i\{\tau_i<\infty\}=1,
\]

и невозвратно, если эта вероятность меньше единицы. Если возвратное
состояние имеет конечное среднее время возвращения, оно положительно
возвратно.

Стационарное распределение - это вероятностный вектор \(\pi\), который
сохраняется переходом:

\[
\pi=\pi P.
\]

Если цепь стартует из \(\pi\), то все \(X_n\) имеют распределение
\(\pi\). Для конечной неразложимой апериодической цепи стационарное
распределение единственно и

\[
p_{ij}^{(n)}\to\pi_j.
\]

Главные ошибки: считать марковость независимостью, забывать начальное
распределение и путать стационарное распределение со стационарностью
цепи.

\subsection{15. Если осталось 2 минуты перед
экзаменом}\label{ux435ux441ux43bux438-ux43eux441ux442ux430ux43bux43eux441ux44c-2-ux43cux438ux43dux443ux442ux44b-ux43fux435ux440ux435ux434-ux44dux43aux437ux430ux43cux435ux43dux43eux43c}

Запомнить:

\[
P\{X_{n+1}=j\mid X_0=i_0,\ldots,X_n=i\}
=p_{ij}.
\]

Закон траектории:

\[
P\{X_0=i_0,\ldots,X_n=i_n\}
=
\pi_0(i_0)p_{i_0i_1}\cdots p_{i_{n-1}i_n}.
\]

Распределения:

\[
\pi_{n+1}=\pi_nP,
\qquad
\pi_n=\pi_0P^n.
\]

Стационарное распределение:

\[
\pi=\pi P.
\]

Характеристики:

\[
i\to j \Leftrightarrow \exists n:\ p_{ij}^{(n)}>0,
\qquad
d(i)=\gcd\{n:p_{ii}^{(n)}>0\}.
\]

Возвратность:

\[
P_i\{\tau_i<\infty\}=1.
\]

Для конечной неразложимой апериодической цепи:

\[
p_{ij}^{(n)}\to\pi_j.
\]

Главная ловушка: марковская цепь - не независимая последовательность;
зависимость есть, но она проходит через текущее состояние.

\newpage

\end{document}
